\documentclass{tnqjournal}
\NewMarkClass{tnq-handoff}
\fancyhead[LE]{\bfseries\fontsize{7.2}{8}\selectfont
  MARKTOP9811-\TopMark{tnq-handoff}\quad
  MARKFIRST9812-\FirstMark{tnq-handoff}\quad
  MARKLAST9813-\LastMark{tnq-handoff}}
\title{LaTeX Mark-Class Handoff Regression}
\abstracttext{The special first page must update the current LaTeX mark-class structures before control returns to the kernel two-column output routine.}
\keywords{mark classes; running head; output handoff}
\authorstub{\textbf{Mark Tester}\par Chennai, India}
\newcommand\markchunk{A running-head mark belongs to the article stream and follows its source position across the transactional split. The independently composed author lane must not alter article mark state. }
\newcommand\markafterchunk{Ordinary native composition retains the mark regions assembled by the LaTeX kernel for the current page and its two columns. }
\begin{document}
\maketitle
\section{Introduction}
\InsertMark{tnq-handoff}{PAGEONE9801}
MARKSTART9804
\markchunk\markchunk\markchunk\markchunk\markchunk
\markchunk\markchunk\markchunk\markchunk\markchunk
\markchunk\markchunk\markchunk\markchunk\markchunk
\markchunk\markchunk\markchunk\markchunk\markchunk
\markchunk\markchunk\markchunk\markchunk\markchunk
\markchunk\markchunk\markchunk\markchunk\markchunk
MARKTAIL9805 The first paragraph continues on the next page.\par
\InsertMark{tnq-handoff}{PAGETWO9802}
\section{Native continuation}
MARKAFTER9806
\markafterchunk\markafterchunk\markafterchunk\markafterchunk\markafterchunk
\markafterchunk\markafterchunk\markafterchunk\markafterchunk\markafterchunk
\markafterchunk\markafterchunk\markafterchunk\markafterchunk\markafterchunk
\markafterchunk\markafterchunk\markafterchunk\markafterchunk\markafterchunk
\end{document}
